% Generated draft: requires mathematical and source review.
\documentclass[11pt]{article}
\usepackage[margin=25mm]{geometry}
\usepackage{fontspec}
\setmainfont{Latin Modern Roman}
\usepackage{amsmath,amssymb,mathtools,unicode-math}
\setmathfont{Latin Modern Math}
\usepackage{enumitem,needspace,xurl,hyperref}
\setlength{\parindent}{0pt}
\setlength{\parskip}{5pt}
\setlength{\emergencystretch}{3em}
\newcommand{\N}{\mathbb{N}}
\newcommand{\Z}{\mathbb{Z}}
\newcommand{\Q}{\mathbb{Q}}
\newcommand{\R}{\mathbb{R}}
\newcommand{\C}{\mathbb{C}}
\newcommand{\F}{\mathbb{F}}
\newcommand{\A}{\mathbb{A}}
\newcommand{\PP}{\mathbb{P}}
\newcommand{\E}{\mathbb{E}}
\newcommand{\eps}{\varepsilon}
\newcommand{\dd}{\,\mathrm{d}}
\begin{document}
\section*{UnsolvedMath ID: 30000505}
\textbf{Source problem code:} catalogue entry
\subsection{Definitions and mathematical statement}
Let \((\mathcal H_n)_{n\geq0}\), with level maps \(\mathcal H_{n+1}\to\mathcal H_n\), be a Modular Tower in the source's intended sense. Such towers ordinarily arise from Hurwitz moduli spaces built from finite-group, prime, conjugacy-class, and successive Frattini-cover data; the actual admissibility assumptions and choice of Hurwitz space are not reproduced. A component branch is a compatible sequence \(C_n\subseteq\mathcal H_n\) of the intended components, with each level map carrying \(C_{n+1}\) into \(C_n\).

The Main Conjecture for Modular Towers in the record is
\[
\forall\text{ admissible component branches }(C_n)\quad
\forall K\text{ a number field}\quad
\exists N=N((C_n),K)\quad
\forall n\geq N,\qquad C_n(K)=\varnothing.
\]
The threshold may depend on the branch and number field; no common threshold over all branches or all fields is requested. The assertion is eventual absence of points on each whole level component, not merely finiteness of its rational points or nonexistence of an infinite compatible sequence of points. Points at different levels need not be compatible for the statement to apply.

The component and field-of-definition conventions are not specified. If the branch consists of geometric components, \(C_n(K)\) can be read as ambient \(K\)-rational points lying on that component, with the precise descent convention to be checked. The relevant open Hurwitz spaces must also be distinguished from their compactifications: adding cusp or degenerate-cover boundary points would change the rational-point statement. The source does not explicitly give the prime, group assumptions, branch-point number, conjugacy classes, inner versus absolute moduli convention, or reduced versus unreduced space. A general inverse system of varieties is not substituted for a Modular Tower satisfying its source restrictions.

The supplied literature assessment checked 2026-08-21 reports the arbitrary-branch, arbitrary-number-field assertion as conjectural and says the cited report gives no proof. The named Main Conjecture is retained without importing hypotheses or claimed special cases from an uncited construction.
\subsection{Short English statement}
Along each admissible Modular Tower component branch, every fixed number field should lose all rational points from sufficiently high levels onward. The level bound may depend on the branch and field, and the intended open moduli and tower hypotheses must be recovered.
\subsection{Sources}
\begin{itemize}
\item Local UnsolvedMath registry, permanent ID 30000505, catalogue entry, \emph{Main Conjecture for Modular Towers}.
\item \emph{Pro-\(p\) Extensions of Global Fields and pro-\(p\) Groups} (2007), Oberwolfach Reports, pp. 1463--1536, DOI \url{https://doi.org/10.4171/owr/2006/25}.
\item Supplied literature triage checked 2026-08-21 cites only the original report; no separate tower-definition reference, exact construction, or proof is supplied.
\end{itemize}
\end{document}
